\section{Introduction}

Bayesian tabulation audits answer a different question from frequentist
risk-limiting audits. In the Rivest-Shen tradition, the audit treats the
reported tabulation and the observed evidence as inputs to a posterior model
and reports the probability that a full hand count would overturn the reported
outcome. That number is useful for explanation and comparison, but it is
conditional on the declared prior and likelihood.

RCOUNT therefore records the Bayesian surface as an analytic boundary. The
verifier does not accept an exported posterior value as a risk limit, and it
does not convert posterior upset probability into frequentist certification.
Instead, it checks that a record naming the Bayesian tabulation method carries
the fields needed to describe the analysis: prior identity, likelihood identity,
posterior winner probability, posterior upset risk, and any declared simulation
or calibration metadata.

This paper covers that implemented V.20 slice. It claims only that the Bayesian
audit record is well formed and kept separate from risk-limiting replay. It
does not claim that RCOUNT currently recomputes a posterior distribution or that
a Bayesian upset probability is itself an RLA risk limit.
